% Reescrever ambos os resumos ao terminar o ensaio; manter equivalencia.
% O tamanho e a quantidade de palavras-chave abaixo sao escolhas editaveis.
\begin{orientacao}{Resumos provisórios}
As duas versões abaixo enunciam o projeto, sem apresentar resultados como
concluídos. Ajuste-as depois de escrever as três partes, sobretudo a
definição de Cecinepá. Para manter a avaliação anônima, examine também
referências autobiográficas, agradecimentos e a descrição desse espaço.
\end{orientacao}

\Needspace{5\baselineskip}
\noindent\textbf{Abstract}\par
\noindent
This essay approaches contemporary systemic crises through Anthropocene
talk: the practices through which a changing world is named, described,
and made available for collective action. A fictional sciences bar runs
alongside the argument, following a Purifier whose efforts to conceal
mediations progressively alter the conditions of speech. Drawing on Bruno
Latour and the controversies surrounding the Anthropocene, the essay
proposes quotation marks as a sign of situated borrowing and translation.
It develops this proposal through Science-bar and Cecinepá, bringing
Latour's common world into conversation with Manuela Carneiro da Cunha's
account of culture in quotation marks, Michel Foucault's heterotopias,
Marcel Mauss's total social fact, Ailton Krenak's ancestral futures, and
Leda Maria Martins's spiral time. The guiding question concerns how to
articulate differences and dependencies while composing conditions for
living together.

\medskip
\noindent\textbf{Keywords:} Anthropocene; quotation marks; Science-bar;
ancestral futures; common world.

\begin{otherlanguage}{portuguese}
\Needspace{5\baselineskip}
\bigskip
\noindent\textbf{Resumo}\par
\noindent
Este ensaio aborda as crises sistêmicas contemporâneas a partir do
\textit{Anthropocene talk}: as práticas pelas quais um mundo em
transformação é nomeado, descrito e tornado disponível à ação coletiva.
Um bar de ciências ficcional acompanha o argumento, seguindo um
Purificador cujos esforços para ocultar mediações alteram progressivamente
as condições da fala. Em diálogo com Bruno Latour e com as controvérsias
em torno do Antropoceno, o ensaio propõe as aspas como sinal de empréstimo
situado e tradução. Desenvolve essa proposta por meio de Science-bar e
Cecinepá, aproximando o mundo comum de Latour da cultura com aspas de
Manuela Carneiro da Cunha, das heterotopias de Michel Foucault, do fato
social total de Marcel Mauss, dos futuros ancestrais de Ailton Krenak e do
tempo espiralar de Leda Maria Martins. A pergunta orientadora concerne a
como articular diferenças e dependências enquanto se compõem condições
de vida em comum.

\medskip
\noindent\textbf{Palavras-chave:} Antropoceno; aspas; Science-bar;
futuros ancestrais; mundo comum.
\end{otherlanguage}
